\chapter{Introduction}
\label{chap:One}

\section{Project Overview}
This thesis addresses the fundamental disconnect between how mathematics is taught to humans and how it is verified by machines. For centuries, mathematical concepts have been communicated using natural language—a medium that, while exceptionally expressive, is inherently ambiguous and prone to implicit assumptions \cite{rosen2019}. In contrast, formal verification systems require rigorous, unambiguous, and structurally flawless logical syntax. This project, highly inspired by contemporary breakthroughs such as Google DeepMind's AlphaProof \cite{alphaproof2024, polu2025} and Harmonic's Aristotle \cite{hubert2025}, aims to bridge the gap between ``Informal English'' and ``Formal Logic''. By leveraging the generative capabilities of Large Language Models (LLMs) alongside the rigorous verification engine of a theorem prover, this thesis proposes an automated continuous pipeline to translate undergraduate-level mathematics into the Lean 4 ecosystem.

\section{Motivation and Significance}
The reliance on human intuition in textual mathematical communication often obscures hidden logical gaps and unspoken axioms. The traditional process of auto-formalizing mathematics is notoriously labor-intensive, often requiring weeks or even months of specialized expert effort to fully translate a single mathematical paper from LaTeX into a formal language \cite{formal_history}. Historically, even highly celebrated mathematicians, such as Fields Medalist Terence Tao, have discovered minor implicit errors or skipped steps in their published mathematical proofs when those proofs were forcibly subjected to formal computer-checked translation \cite{hubert2025}.

With the rapid advent of powerful Large Language Models (LLMs), there is now an unprecedented opportunity to automate this translation process, theoretically eliminating the bottleneck of manual expert labor \cite{llm_math}. However, LLMs are famously prone to mathematical ``hallucinations''—they frequently generate highly plausible but logically flawed, incorrect, or entirely fictitious mathematical statements \cite{llm_hallucination}. Therefore, pairing the creative pattern-matching capability of an LLM with a strict, symbolic, cryptographic compiler like Lean 4 ensures a flawless end-to-end system: the LLM provides the intuition and the raw code, while Lean 4 acts as the ultimate, unyielding arbiter of mathematical truth \cite{neural_formal}.

\section{Objectives}
The core objective of this research project is to conceptualize, design, and build a ``Truth Scanner'' pipeline. This pipeline must be capable of automatically translating foundational definitions, theorems, and proofs from a highly-regarded academic textbook—specifically Kenneth Rosen's widely used \textit{Discrete Mathematics and Its Applications} \cite{rosen2019}—into functional, compiler-verified Lean 4 logic \cite{lean4_doc}. 

The specific sub-objectives of this thesis include:
\begin{enumerate}
    \item \textbf{Decomposition:} Break down complex English mathematical definitions into atomic logical lemmas that an AI model can consistently parse accurately.
    \item \textbf{Ground Truth Formalization:} Manually construct a baseline ``Rosetta Stone'' repository of valid definitions in Lean 4 to serve as a high-quality benchmark for evaluating the LLM's automated output.
    \item \textbf{Prompt Engineering API:} Develop targeted prompt-engineering strategies that guide commercial LLMs (e.g., OpenAI's GPT-4, Google's Gemini) to seamlessly output structurally valid Lean 4 syntax \cite{llm_codegen}.
    \item \textbf{Autonomous Repair Loop:} Implement a sophisticated Python-based self-correcting system where the Lean 4 compiler's output and error traces are intercepted and autonomously fed back into the LLM, enabling the AI to ``debug'' its own mathematical logic until the compiler passes.
\end{enumerate}

\section{Thesis Organization}
This thesis is structured to provide a clear narrative of the automation pipeline. Chapter \ref{chap:Two} reviews the fundamental background on formal theorem proving, the history of auto-formalization, and the architectural mechanics of Lean 4. Chapter \ref{chap:Three} details the system methodology, providing visual architectures of the proposed repair loops and Truth Scanner pipeline. Chapter \ref{chap:Four} evaluates the manually formalized framework and discusses early extraction results. Finally, Chapter \ref{chap:Five} concludes the thesis by reflecting on the intersection of Artificial Intelligence and rigorous Mathematics, while outlining future expansion possibilities.